\documentclass[aspectratio=169,10pt]{beamer}
\usetheme[numbering=fraction,progressbar=frametitle]{metropolis}
\usepackage{fontspec}
\setsansfont{Latin Modern Sans}
\usepackage{amsmath,amssymb,mathtools,graphicx,booktabs}
\usepackage{polyglossia}
\setdefaultlanguage{english}
\setbeamertemplate{navigation symbols}{}
\setbeamerfont{frametitle}{size=\large}
\setbeamerfont{normal text}{size=\normalsize}
\setlength{\parskip}{.45em}
\hypersetup{pdfauthor={},pdfsubject={Lecture slides},pdfkeywords={},pdftitle={The reciprocal function}}
\providecommand{\N}{\mathbb{N}}
\providecommand{\Z}{\mathbb{Z}}
\providecommand{\D}{\mathbb{D}}
\providecommand{\Q}{\mathbb{Q}}
\providecommand{\R}{\mathbb{R}}
\providecommand{\C}{\mathbb{C}}
\providecommand{\vect}[1]{\overrightarrow{#1}}
\providecommand{\vi}{\vec{\imath}}
\providecommand{\vj}{\vec{\jmath}}
\providecommand{\vk}{\vec{k}}
\providecommand{\norm}[1]{\left\lVert #1\right\rVert}
\providecommand{\vc}[2]{\begin{pmatrix}#1\\#2\end{pmatrix}}
\providecommand{\deter}[2]{\det\left(#1,#2\right)}
\providecommand{\abs}[1]{\left\lvert #1\right\rvert}
\providecommand{\dd}{\,\mathrm{d}}
\providecommand{\e}{\mathrm{e}}
\providecommand{\Card}{\operatorname{Card}}
\providecommand{\E}{\operatorname{E}}
\providecommand{\V}{\operatorname{V}}
\providecommand{\degre}{^{\circ}}
\providecommand{\un}[1]{\,\mathrm{#1}}
\providecommand{\sci}[2]{#1\times10^{#2}}
\providecommand{\pgcd}{\operatorname{PGCD}}
\providecommand{\ppcm}{\operatorname{PPCM}}
\title{The reciprocal function}
\subtitle{Lecture and worked examples}
\date{}
\author{}
\begin{document}
\begin{frame}[plain]
\vfill
{\usebeamerfont{title}\inserttitle\par}
\medskip
{\usebeamerfont{subtitle}\insertsubtitle\par}
\medskip
{\usebeamercolor[fg]{progress bar}\rule{\textwidth}{.5pt}}
\vfill
\end{frame}
\begin{frame}{Reference notes}
\(f'(x)=-\dfrac1{x^2}<0\) for every \(x\neq0\). The reciprocal function is strictly decreasing on each of the intervals \(]-\infty,0[\) and \(]0,+\infty[\). The lines \(x=0\) and \(y=0\) are asymptotes.
\end{frame}
\begin{frame}{Reference notes}
One cannot write that the reciprocal function is decreasing “on \(\R\setminus\{0\}\)” without specifying the two intervals. For example \(-1<1\) but \(f(-1)=-1<1=f(1)\), which would contradict the global definition of decreasing.
\end{frame}
\begin{frame}{Reference notes}
On \(]0,+\infty[\), \(g(x)=x+\dfrac4x\) satisfies \(g'(x)=1-\dfrac4{x^2}=\dfrac{(x-2)(x+2)}{x^2}\). The denominator is positive; \(g'\) is negative on \(]0,2[\) then positive on \(]2,+\infty[\). The function has a minimum \(g(2)=4\).
\end{frame}
\begin{frame}{Definition and graph of the reciprocal function}
\begin{itemize}
\item Find the domain.
\item Draw the curve.
\item Use symmetries.

\end{itemize}
\end{frame}
\begin{frame}{Definition and graph of the reciprocal function}
\begin{columns}[T,onlytextwidth]
\begin{column}{.57\textwidth}
Calculate \(\dfrac1x\) for \(x=-4,-2,-1,-\dfrac12,\dfrac12,1,2,4\). Plot the points and describe the symmetries. Why is \(x=0\) missing?
center center
\end{column}
\begin{column}{.4\textwidth}
\centering
\includegraphics[width=\linewidth,height=4.5cm,keepaspectratio]{figures/stmg-04-2-en.pdf}
\end{column}
\end{columns}
\end{frame}
\begin{frame}{Definition and graph of the reciprocal function}
The reciprocal function is \(f(x)=\dfrac1x\) on \(\R\setminus\{0\}\). Its curve is a hyperbola symmetric about the origin.
\end{frame}
\begin{frame}{Definition and graph of the reciprocal function}
Revisit the data from the investigation and write up the solution in order: definition of the objects, property used, exact or approximate calculation, then a concluding sentence. The model answer appears in the teacher edition.
\end{frame}
\begin{frame}{Definition and graph of the reciprocal function}
\textbf{1 h session:} 7 min of fluency; 10 min of pooling and written record; 28 min of essential exercises; 5 min of correction; 3 min of exit ticket.

\textbf{2 h session:} add 20 min of investigation, 20 min of formalisation, a second independent exercise and 10 min of comparative correction.

\textbf{Differentiation:} give struggling students a four-box sheet (data, model, calculation, conclusion). Ask quick students to produce a counterexample, a justification or a second method.
\end{frame}
\begin{frame}{Behaviour at the boundaries and asymptotes}
\begin{itemize}
\item Describe the behaviour at the boundaries.
\item Identify the asymptotes.
\item Read a graph.

\end{itemize}
\end{frame}
\begin{frame}{Behaviour at the boundaries and asymptotes}
\begin{columns}[T,onlytextwidth]
\begin{column}{.57\textwidth}
Using a table of values near \(0\) and for large values of \(|x|\), describe \(\dfrac1x\). Translate each behaviour into an asymptote and a sentence.
center center
\end{column}
\begin{column}{.4\textwidth}
\centering
\includegraphics[width=\linewidth,height=4.5cm,keepaspectratio]{figures/stmg-04-3-en.pdf}
\end{column}
\end{columns}
\end{frame}
\begin{frame}{Behaviour at the boundaries and asymptotes}
For \(f(x)=\dfrac1x\), \(x=0\) is a vertical asymptote and \(y=0\) a horizontal asymptote. The branches approach them without crossing them.
\end{frame}
\begin{frame}{Behaviour at the boundaries and asymptotes}
Revisit the data from the investigation and write up the solution in order: definition of the objects, property used, exact or approximate calculation, then a concluding sentence. The model answer appears in the teacher edition.
\end{frame}
\begin{frame}{Behaviour at the boundaries and asymptotes}
\textbf{1 h session:} 7 min of fluency; 10 min of pooling and written record; 28 min of essential exercises; 5 min of correction; 3 min of exit ticket.

\textbf{2 h session:} add 20 min of investigation, 20 min of formalisation, a second independent exercise and 10 min of comparative correction.

\textbf{Differentiation:} give struggling students a four-box sheet (data, model, calculation, conclusion). Ask quick students to produce a counterexample, a justification or a second method.
\end{frame}
\begin{frame}{Derivative and direction of variation}
\begin{itemize}
\item Differentiate the reciprocal function.
\item Justify the variations.
\item Draw up a table.

\end{itemize}
\end{frame}
\begin{frame}{Derivative and direction of variation}
Starting from the difference quotient of \(f(x)=\dfrac1x\), recover \(f'(x)=-\dfrac1{x^2}\). Study its sign and check consistency with the graph.
\end{frame}
\begin{frame}{Derivative and direction of variation}
For \(x\neq0\), \(f'(x)=-\dfrac1{x^2}<0\). The reciprocal function decreases on each of the intervals \(]-\infty,0[\) and \(]0,+\infty[\).
\end{frame}
\begin{frame}{Derivative and direction of variation}
Revisit the data from the investigation and write up the solution in order: definition of the objects, property used, exact or approximate calculation, then a concluding sentence. The model answer appears in the teacher edition.
\end{frame}
\begin{frame}{Derivative and direction of variation}
\textbf{1 h session:} 7 min of fluency; 10 min of pooling and written record; 28 min of essential exercises; 5 min of correction; 3 min of exit ticket.

\textbf{2 h session:} add 20 min of investigation, 20 min of formalisation, a second independent exercise and 10 min of comparative correction.

\textbf{Differentiation:} give struggling students a four-box sheet (data, model, calculation, conclusion). Ask quick students to produce a counterexample, a justification or a second method.
\end{frame}
\begin{frame}{Combinations with polynomial functions}
\begin{itemize}
\item Differentiate a combination of functions.
\item Simplify an expression.
\item Organise a function study.

\end{itemize}
\end{frame}
\begin{frame}{Combinations with polynomial functions}
A company produces \(x>0\) hundred items and its average cost is \(C(x)=2x+30+\dfrac{200}{x}\). Calculate \(C'(x)\), factorise its numerator and find the production level that minimises the average cost.
\end{frame}
\begin{frame}{Combinations with polynomial functions}
To study \(P(x)+\dfrac{k}{x}\), find the domain, calculate the derivative, put over a common denominator, study its sign, then draw up the variations.
\end{frame}
\begin{frame}{Combinations with polynomial functions}
Revisit the data from the investigation and write up the solution in order: definition of the objects, property used, exact or approximate calculation, then a concluding sentence. The model answer appears in the teacher edition.
\end{frame}
\begin{frame}{Combinations with polynomial functions}
\textbf{1 h session:} 7 min of fluency; 10 min of pooling and written record; 28 min of essential exercises; 5 min of correction; 3 min of exit ticket.

\textbf{2 h session:} add 20 min of investigation, 20 min of formalisation, a second independent exercise and 10 min of comparative correction.

\textbf{Differentiation:} give struggling students a four-box sheet (data, model, calculation, conclusion). Ask quick students to produce a counterexample, a justification or a second method.
\end{frame}
\begin{frame}{Sign and factorisation}
\begin{itemize}
\item Factorise.
\item Build a sign table.
\item Respect excluded values.

\end{itemize}
\end{frame}
\begin{frame}{Sign and factorisation}
Study the sign of \(\dfrac{(x-2)(x+3)}x\). Build the table with the three critical values, distinguish a zero from an excluded value, then solve the inequality \(\dfrac{(x-2)(x+3)}x\geq0\).
\end{frame}
\begin{frame}{Sign and factorisation}
In a quotient, the sign depends on the numerator and the denominator. Excluded values are left out of the table, even if a factor formally cancels.
\end{frame}
\begin{frame}{Sign and factorisation}
Revisit the data from the investigation and write up the solution in order: definition of the objects, property used, exact or approximate calculation, then a concluding sentence. The model answer appears in the teacher edition.
\end{frame}
\begin{frame}{Sign and factorisation}
\textbf{1 h session:} 7 min of fluency; 10 min of pooling and written record; 28 min of essential exercises; 5 min of correction; 3 min of exit ticket.

\textbf{2 h session:} add 20 min of investigation, 20 min of formalisation, a second independent exercise and 10 min of comparative correction.

\textbf{Differentiation:} give struggling students a four-box sheet (data, model, calculation, conclusion). Ask quick students to produce a counterexample, a justification or a second method.
\end{frame}
\begin{frame}{Derivatives of polynomials of degree at most 3}
\begin{itemize}
\item Differentiate a polynomial.
\item Study a derivative.
\item Deduce the variations.

\end{itemize}
\end{frame}
\begin{frame}{Derivatives of polynomials of degree at most 3}
For \(f(x)=2x^3-3x^2-12x+5\), calculate \(f'(x)\), factorise, study its sign and deduce the variations of \(f\).
\end{frame}
\begin{frame}{Derivatives of polynomials of degree at most 3}
\((ax^3+bx^2+cx+d)'=3ax^2+2bx+c\). The sign of the derivative gives the direction of variation of the function.
\end{frame}
\begin{frame}{Derivatives of polynomials of degree at most 3}
Revisit the data from the investigation and write up the solution in order: definition of the objects, property used, exact or approximate calculation, then a concluding sentence. The model answer appears in the teacher edition.
\end{frame}
\begin{frame}{Derivatives of polynomials of degree at most 3}
\textbf{1 h session:} 7 min of fluency; 10 min of pooling and written record; 28 min of essential exercises; 5 min of correction; 3 min of exit ticket.

\textbf{2 h session:} add 20 min of investigation, 20 min of formalisation, a second independent exercise and 10 min of comparative correction.

\textbf{Differentiation:} give struggling students a four-box sheet (data, model, calculation, conclusion). Ask quick students to produce a counterexample, a justification or a second method.
\end{frame}
\begin{frame}{Full study and graph}
\begin{itemize}
\item Carry out a full study.
\item Draw a consistent curve.
\item Interpret an extremum.

\end{itemize}
\end{frame}
\begin{frame}{Full study and graph}
\begin{columns}[T,onlytextwidth]
\begin{column}{.57\textwidth}
Study \(g(x)=x+\dfrac9x\) completely on \(]0,+\infty[\): domain, derivative, sign, variations, minimum, behaviour at the boundaries and sketch.
center center
\end{column}
\begin{column}{.4\textwidth}
\centering
\includegraphics[width=\linewidth,height=4.5cm,keepaspectratio]{figures/stmg-04-4-en.pdf}
\end{column}
\end{columns}
\end{frame}
\begin{frame}{Full study and graph}
A full study includes: domain, derivative, sign, variations, limits or behaviour at the boundaries, key points and a consistent graph.
\end{frame}
\begin{frame}{Full study and graph}
Revisit the data from the investigation and write up the solution in order: definition of the objects, property used, exact or approximate calculation, then a concluding sentence. The model answer appears in the teacher edition.
\end{frame}
\begin{frame}{Full study and graph}
\textbf{1 h session:} 7 min of fluency; 10 min of pooling and written record; 28 min of essential exercises; 5 min of correction; 3 min of exit ticket.

\textbf{2 h session:} add 20 min of investigation, 20 min of formalisation, a second independent exercise and 10 min of comparative correction.

\textbf{Differentiation:} give struggling students a four-box sheet (data, model, calculation, conclusion). Ask quick students to produce a counterexample, a justification or a second method.
\end{frame}
\begin{frame}{Review and assessment}
\begin{itemize}
\item Use domain, derivative and sign.
\item Write up a study.
\item Correct work methodically.

\end{itemize}
\end{frame}
\begin{frame}{Review and assessment}
Correct a deliberately flawed study of \(h(x)=x+\dfrac4x\): stated domain \(\R\), derivative \(1+\dfrac4{x^2}\), table with no break at \(0\). Rewrite a complete solution.
\end{frame}
\begin{frame}{Review and assessment}
Check the four items separately: domain, expression of the derivative, sign of the derivative, variations. An error at one step must not hide the following ones.
\end{frame}
\begin{frame}{Review and assessment}
Revisit the data from the investigation and write up the solution in order: definition of the objects, property used, exact or approximate calculation, then a concluding sentence. The model answer appears in the teacher edition.
\end{frame}
\begin{frame}{Review and assessment}
\textbf{1 h session:} 7 min of fluency; 10 min of pooling and written record; 28 min of essential exercises; 5 min of correction; 3 min of exit ticket.

\textbf{2 h session:} add 20 min of investigation, 20 min of formalisation, a second independent exercise and 10 min of comparative correction.

\textbf{Differentiation:} give struggling students a four-box sheet (data, model, calculation, conclusion). Ask quick students to produce a counterexample, a justification or a second method.
\end{frame}
\end{document}
